\documentclass[11pt]{article}

\usepackage[utf8]{inputenc}
\usepackage[T1]{fontenc}
\usepackage{amsmath, amssymb, amsthm}
\usepackage{mathtools}
\usepackage[margin=1in]{geometry}
\usepackage{hyperref}

\theoremstyle{plain}
\newtheorem{theorem}{Theorem}[section]
\newtheorem{lemma}[theorem]{Lemma}
\newtheorem{corollary}[theorem]{Corollary}

\theoremstyle{definition}
\newtheorem{definition}[theorem]{Definition}
\newtheorem*{remark}{Remark}
\newtheorem*{notation}{Notation}
\newtheorem*{assertion}{Assertion}
\newtheorem*{claim}{Claim}

\newcommand{\R}{\mathbb{R}}
\newcommand{\C}{\mathbb{C}}
\newcommand{\Z}{\mathbb{Z}}
\newcommand{\Q}{\mathbb{Q}}
\newcommand{\nbd}{neighborhood}
\DeclareMathOperator{\ord}{ord}
\DeclareMathOperator{\res}{res}
\DeclareMathOperator{\disc}{disc}
\DeclareMathOperator{\coeff}{coeff}

\title{Generalized Lifting and Projection Theorems\\for Cylindrical Algebraic Decomposition}
\author{Generalization of McCallum (1985), Chapter 3}
\date{Working document --- \today}

\begin{document}

\maketitle

\begin{abstract}
We state and prove generalizations of McCallum's Lifting Theorem (3.2.1)
and Projection Theorem (3.2.3).  The generalized versions replace the
discriminant $D(x) = \disc(f)$ and specific resultants $\res(F_i, F_j)$
by arbitrary nonzero elements of the corresponding elimination ideals
$\langle f, \partial f / \partial x_r \rangle \cap \R[x]$ and
$\langle F_i, F_j \rangle \cap \R[x]$.  The proofs follow the structure of
McCallum's original argument with two new ingredients: (1) the
\emph{order additivity lemma} for holomorphic functions, and
(2) the fact that the elimination ideal
$\langle h, h' \rangle \cap \mathcal{O}(\Delta)$ is principal (generated by the
discriminant) for a Weierstrass polynomial~$h$.

Sections that are unchanged from the original thesis are marked
\textbf{[UNCHANGED]} and only summarized.
Sections that require modification are marked \textbf{[CHANGED]} and
written out in full.
\end{abstract}

\tableofcontents

%% =========================================================================
\section{Setup and notation}
\label{sec:setup}
%% =========================================================================

\textbf{[UNCHANGED]}
All notation, definitions, and background from Chapter~2 of the thesis
carry over without modification.  In particular:
\begin{itemize}
\item $K = \R$ or $\C$; analytic/holomorphic functions; power series
  (Sec.~2.1).
\item Submanifolds, coordinate systems, charts (Sec.~2.2, Theorems 2.2.1--2.2.4).
\item Order $\ord_p f$; order-invariance; Theorem~2.3.1 (order under
  composition); root continuity principle (Theorem~2.3.2).
\item Resultant, discriminant, Theorem~2.3.3 (discriminant of a product).
\item Weierstrass preparation theorem (Theorem~2.1.10).
\item Zariski's 1975 theorem (Theorem~4.1.1) and its proof (Chapter~4).
\end{itemize}

Throughout, $x = (x_1, \ldots, x_{r-1})$ and $(x, x_r)$ denotes an
$r$-tuple.  A polynomial $f(x, x_r) \in \R[x, x_r]$ is viewed as
univariate in~$x_r$ over $\R[x]$.

%% =========================================================================
\section{The generalized projection map \texorpdfstring{$P'(A)$}{P'(A)}}
\label{sec:projection_map}
%% =========================================================================

\textbf{[CHANGED]}

\begin{definition}[Squarefree basis]
A finite set $A$ of polynomials in $\R[x, x_r]$ is a \emph{squarefree
basis} if every element has positive degree in~$x_r$, every element is
squarefree (in $\R[x, x_r]$), and the elements are pairwise coprime.
\end{definition}

\begin{definition}[Generalized reduced projection]
\label{def:gen_projection}
Let $A$ be a squarefree basis.  The \emph{generalized reduced projection}
$P'(A)$ is a set of polynomials in $\R[x]$ of the form
\[
  P'(A) \;=\; \coeff(A) \;\cup\; \mathrm{elim\_self}(A)
              \;\cup\; \mathrm{elim\_pair}(A),
\]
where:
\begin{itemize}
\item $\coeff(A) = \{ c \in \R[x] : c$ is a nonzero coefficient of some
  $F \in A$ viewed as a polynomial in~$x_r \}$.
\item $\mathrm{elim\_self}(A) = \{ d_F : F \in A,\; \deg(F) \ge 2 \}$,
  where each $d_F$ is a \emph{chosen} nonzero element of
  $\langle F, F' \rangle \cap \R[x]$.
  (Here $F' = \partial F / \partial x_r$.)
\item $\mathrm{elim\_pair}(A) = \{ r_{F,G} : F, G \in A,\; F \ne G,\;
  \deg(F) \ge 1,\; \deg(G) \ge 1 \}$, where each $r_{F,G}$ is a
  \emph{chosen} nonzero element of $\langle F, G \rangle \cap \R[x]$.
\end{itemize}
\end{definition}

\begin{remark}
Taking $d_F = \disc(F)$ and $r_{F,G} = \res(F, G)$ recovers the
original projection $P(A) = \coeff(A) \cup \disc(A) \cup \res(A)$.
The discriminant and resultant are particular nonzero elements of the
respective elimination ideals (see Lemma~\ref{lem:discr_in_elim}).
\end{remark}

\begin{lemma}
\label{lem:discr_in_elim}
Let $F \in \R[x, x_r]$ be squarefree of degree $\ge 2$ in~$x_r$.
Then $\disc(F) \ne 0$ and $\disc(F) \in \langle F, F' \rangle \cap \R[x]$
(up to a unit factor involving the leading coefficient).

Similarly, for $F \ne G$ coprime in $\R[x, x_r]$,
$\res(F, G) \ne 0$ and
$\res(F, G) \in \langle F, G \rangle \cap \R[x]$.
\end{lemma}
\begin{proof}
Standard properties of resultants and discriminants:
$\res(F, G) \in \langle F, G \rangle$ follows from the
Sylvester matrix identity, and $\disc(F) = \pm \res(F, F') / \mathrm{lc}(F)$.
\end{proof}

%% =========================================================================
\section{Prerequisite: order-invariance of products (Lemma~3.2.2)}
\label{sec:lemma_322}
%% =========================================================================

\textbf{[UNCHANGED]}

\begin{lemma}[Lemma 3.2.2]
\label{lem:order_mul}
Let $K = \R$ or $\C$.  Let $F_1, \ldots, F_n$ be nonconstant analytic
functions defined in an open subset $U$ of~$K^r$, and let
$f = F_1 \cdots F_n$.  Let $S$ be a connected subset of~$U$.  Then $f$
is order-invariant in~$S$ if and only if each $F_i$ is order-invariant
in~$S$.
\end{lemma}
\begin{proof}
By additivity of the vanishing order:
$\ord_p(f) = \sum_i \ord_p(F_i)$.
\end{proof}

%% =========================================================================
\section{New prerequisite: order additivity lemma}
\label{sec:order_additivity}
%% =========================================================================

\textbf{[NEW]}

\begin{lemma}[Order additivity lemma]
\label{lem:order_additivity}
Let $U \subseteq \C^s$ be a connected open set.  Let $f, g : U \to \C$
be holomorphic functions, both not identically zero.  If $f \cdot g$
has constant vanishing order on~$U$ (i.e., $\ord_z(fg)$ is the same for
all $z \in U$), then $f$ and $g$ each have constant vanishing order
on~$U$.
\end{lemma}

\begin{proof}
Let $c = \ord_z(fg)$ (constant) for all $z \in U$.
Let $a = \min_{z \in U} \ord_z(f)$ and $b = \min_{z \in U} \ord_z(g)$.

Since $f$ is holomorphic and not identically zero, the set
$\{ z \in U : \ord_z(f) \ge k \}$ is a closed analytic subset of~$U$ for
each $k$.  The order achieves its minimum $a$ on a dense open subset
$U_1 \subseteq U$ (the complement of a proper analytic subset).
Similarly, $\ord_z(g)$ achieves its minimum $b$ on a dense open
$U_2 \subseteq U$.

On $U_1 \cap U_2$ (which is dense and open, since $U$ is connected and
each $U_i$ is the complement of a proper analytic subset), we have
$a + b = c$.

For an arbitrary point $z_0 \in U$:
\begin{itemize}
\item $\ord_{z_0}(f) \ge a$ \quad (since $a$ is the minimum),
\item $\ord_{z_0}(g) \ge b$ \quad (since $b$ is the minimum),
\item $\ord_{z_0}(f) + \ord_{z_0}(g) = c = a + b$.
\end{itemize}
The third equality together with the two inequalities forces
$\ord_{z_0}(f) = a$ and $\ord_{z_0}(g) = b$.
\end{proof}

\begin{remark}
This lemma generalizes to any connected analytic manifold, and also to the
real-analytic setting (by complexification).  The key input is that for a
non-identically-zero holomorphic function on a connected open set, the
vanishing order achieves a finite minimum on a dense open subset and can
only jump \emph{up} on a proper analytic subset.
\end{remark}

%% =========================================================================
\section{New prerequisite: elimination ideal for Weierstrass polynomials}
\label{sec:elim_weierstrass}
%% =========================================================================

\textbf{[NEW]}

\begin{lemma}[Elimination ideal for monic polynomials]
\label{lem:elim_monic}
Let $R$ be an integral domain, and let $h \in R[z_r]$ be monic of
degree~$m$.  Let $g \in R[z_r]$ have $\deg(g) \le m - 1$.  Then
\[
  \langle h, g \rangle \cap R \;=\; \bigl(\mathrm{Res}_{z_r}(h, g)\bigr).
\]
In particular, every element of $\langle h, g \rangle \cap R$ is divisible
by $\mathrm{Res}_{z_r}(h, g)$ in~$R$.
\end{lemma}
\begin{proof}
Standard result from commutative algebra: for $h$ monic, the
resultant $\res_{z_r}(h, g)$ generates the elimination ideal.
(See, e.g., Lang \emph{Algebra}, Ch.~IV, or Cox--Little--O'Shea
\emph{Ideals, Varieties, and Algorithms}, Ch.~3.)
\end{proof}

\begin{corollary}
\label{cor:elim_weierstrass}
Let $h(z, z_r) \in \mathcal{O}(\Delta)[z_r]$ be a Weierstrass polynomial
of degree~$m$ in a polydisc $\Delta \subseteq \C^{n-1}$, and let
$h' = \partial h / \partial z_r$.  Then
\[
  \langle h, h' \rangle \cap \mathcal{O}(\Delta)
  \;=\; \bigl(\disc(h)\bigr)
\]
as an ideal in $\mathcal{O}(\Delta)$.  In particular, any
$P \in \langle h, h' \rangle \cap \mathcal{O}(\Delta)$ can be written as
$P = \disc(h) \cdot Q$ for some $Q \in \mathcal{O}(\Delta)$.
\end{corollary}
\begin{proof}
Since $h$ is monic of degree~$m$ and $h'$ has degree $m - 1$ with
leading coefficient~$m$ (a unit in $\mathcal{O}(\Delta)$ when
$\mathrm{char} = 0$), Lemma~\ref{lem:elim_monic} gives
$\langle h, h' \rangle \cap \mathcal{O}(\Delta) = (\res_{z_r}(h, h'))$.
Since $h$ is monic,
$\res_{z_r}(h, h') = (-1)^{m(m-1)/2} \disc(h)$, so
$(\res_{z_r}(h, h')) = (\disc(h))$.
\end{proof}

%% =========================================================================
\section{The Generalized Lifting Theorem (Theorem~3.2.1')}
\label{sec:lifting}
%% =========================================================================

\textbf{[CHANGED: statement and proof]}

\begin{theorem}[Generalized Lifting Theorem, Theorem 3.2.1']
\label{thm:gen_lifting}
Let $f(x, x_r) \in \R[x, x_r]$ be squarefree, of positive degree
in~$x_r$.  Let $S$ be a connected submanifold of\/ $\R^{r-1}$ in which
$f$ is degree-invariant and of non-negative degree.  Suppose that some
nonzero
\[
  P \;\in\; \langle f, \partial f / \partial x_r \rangle \cap \R[x]
\]
is order-invariant in~$S$.  Then $f$ is analytically delineable on~$S$
and is order-invariant in each $f$-section over~$S$.
\end{theorem}

\begin{remark}
This generalizes the original Theorem~3.2.1, which takes
$P = D(x) = \disc(f)$.  Since $\disc(f) \in \langle f, f' \rangle \cap \R[x]$
(Lemma~\ref{lem:discr_in_elim}) and $\disc(f) \ne 0$ for squarefree~$f$,
the original theorem is a special case.
\end{remark}

\subsection{Proof structure}
\label{ssec:proof_structure}

\textbf{[UNCHANGED]} The proof follows the same overall structure as the
original.  Let the dimension of~$S$ be~$s$, where $0 \le s \le r-1$.
If $s = 0$ the theorem is trivial.  For $1 \le s \le r-1$, the theorem
follows from the following assertion about ``local delineability'':

\begin{assertion}[Assertion 1]
For each point $p$ of~$S$ there is a neighborhood
$N \subseteq \R^{r-1}$ of~$p$ such that $f$ is analytically delineable on
$S \cap N$ and $f$ is order-invariant in each $f$-section over $S \cap N$.
\end{assertion}

\textbf{[UNCHANGED]}
The deduction of the theorem from Assertion~1 is identical to the
original: by connectedness of~$S$, the number of distinct real roots
and their multiplicities are locally constant, hence globally constant.
The local analytic root functions patch to global root functions on~$S$.

\textbf{[UNCHANGED]}
It remains to establish Assertion~1.  As in the original, we reduce to
the case $\alpha_i = 0$ (by translation in~$x_r$), fix
$m = m_i$, $C = C_i$, and seek a neighborhood $N = N_i$ of~$p$ and an
analytic function $\theta : S \cap N \to (-\epsilon, +\epsilon)$
satisfying (3.3.1) and order-invariance.

We handle two cases.

%% -------------------------------------------------------------------------
\subsection{Case \texorpdfstring{$s = r - 1$}{s = r - 1}
  (S is open)}
\label{ssec:case_open}
%% -------------------------------------------------------------------------

\textbf{[CHANGED --- simplified]}

Since $s = r - 1$, $S$ is an open subset of~$\R^{r-1}$.  Let $U \subseteq N_0$
be a neighborhood of~$p$ with $U \subseteq S$.

Since $P$ is order-invariant on~$S$ and $P \ne 0$, and $S$ is open and
connected, $P$ has constant finite order on~$S$.  On an open connected set,
a holomorphic (indeed polynomial) function with constant vanishing order
must be nowhere-vanishing (order~$0$) unless it is identically zero.
Since $P \ne 0$, we have $P(q) \ne 0$ for every $q \in S$.

Since $P \in \langle f, f' \rangle \cap \R[x]$, there exist
$A, B \in \R[x, x_r]$ such that
\[
  P(x) = A(x, x_r) f(x, x_r) + B(x, x_r) \frac{\partial f}{\partial x_r}(x, x_r).
\]
If $f(p, \alpha) = 0$ and $\frac{\partial f}{\partial x_r}(p, \alpha) = 0$
for some~$\alpha$, then evaluating the identity at $(p, \alpha)$ gives
$P(p) = 0$, contradicting $P(p) \ne 0$.  Therefore $f(p, x_r)$ and
$\frac{\partial f}{\partial x_r}(p, x_r)$ have no common root, i.e., all
roots of $f(p, x_r)$ are simple.

In particular, $m = 1$ (the root $\alpha_i = 0$ is simple), and
$\frac{\partial f}{\partial x_r}(p, 0) \ne 0$.  By the implicit function
theorem (Theorem~2.1.7 with $s = r - 1$), there exist neighborhoods
$N' \subseteq U$ and $W \subseteq (-\epsilon, +\epsilon)$ of~$0$, and an
analytic function $\theta : S \cap N' \to W$ such that for all
$x \in S \cap N'$ and $\alpha \in W$,
\[
  f(x, \alpha) = 0 \quad\text{if and only if}\quad \alpha = \theta(x).
\]
Since $m = 1$, $\frac{\partial f}{\partial x_r}(x, \theta(x)) \ne 0$,
so $\ord_{(x, \theta(x))} f = 1$ for all $x \in S \cap N'$.  Hence $f$
is order-invariant (with constant order~$1$) in the graph of~$\theta$.
This establishes Assertion~2 (and hence Assertion~1) in this case.

\begin{remark}
Compare with the original proof (pp.~23--24 of the thesis), which uses
the relation $(-1)^{m(m-1)/2} R(x) = f_0(x) D(x)$ to conclude
$R(p) \ne 0$.  Our argument is simpler: $P(p) \ne 0$ directly rules out
a common root.
\end{remark}

%% -------------------------------------------------------------------------
\subsection{Case \texorpdfstring{$1 \le s \le r - 2$}{1 <= s <= r - 2}
  (S has positive codimension)}
\label{ssec:case_codim}
%% -------------------------------------------------------------------------

\textbf{[CHANGED --- the factorization step is different; everything
  else is unchanged]}

\subsubsection{Coordinate change (unchanged)}
\label{sssec:coord_change}

By Theorem~2.2.1, choose a coordinate system $\Phi : U \to V$ about $p$
with $U \subseteq N_0$ connected, such that
$S \cap U = \{ x \in U : \phi_{s+1}(x) = 0, \ldots, \phi_{r-1}(x) = 0 \}$.
Let $T$ be the $s$-dimensional linear subspace
$\{ y \in \R^{r-1} : y_{s+1} = \cdots = y_{r-1} = 0 \}$, so the image of
$S \cap U$ under~$\Phi$ is $T \cap V$.

Define $g(y, y_r)$ by transforming $f(x, x_r)$ via~$\Phi$:
$g_j(y) = f_j(\Phi^{-1}(y))$ for $0 \le j \le d$, so
$g(y, y_r) = \sum_{j=0}^{d} g_j(y) y_r^{d-j}$ is a pseudopolynomial
with analytic coefficients.

\subsubsection{Transform of \texorpdfstring{$P$}{P} (changed)}
\label{sssec:transform_P}

\textbf{[CHANGED]}

Define $\widetilde{P}(y) = P(\Phi^{-1}(y))$ for $y \in V$.  Since
$P \in \R[x]$ and $\Phi^{-1}$ is analytic, $\widetilde{P}$ is a
real-analytic function in~$V$, not identically zero (as $P \ne 0$).

Since $P(x) = A(x, x_r) f(x, x_r) + B(x, x_r) f'(x, x_r)$ in
$\R[x, x_r]$, substituting $x = \Phi^{-1}(y)$ gives
\[
  \widetilde{P}(y) = \widetilde{A}(y, y_r) \, g(y, y_r)
    + \widetilde{B}(y, y_r) \, \frac{\partial g}{\partial y_r}(y, y_r)
\]
where $\widetilde{A}(y, y_r) = A(\Phi^{-1}(y), y_r)$ and
$\widetilde{B}(y, y_r) = B(\Phi^{-1}(y), y_r)$ are analytic (not
polynomial) in~$(y, y_r)$.  Hence
\[
  \widetilde{P} \;\in\;
    \langle g, \partial g / \partial y_r \rangle \cap \mathcal{O}(V),
\]
where $\mathcal{O}(V)$ denotes the ring of real-analytic functions of~$y$
in~$V$.

Since $P$ is order-invariant in~$S$ and $\ord_{G(p)} f \le \ord_p f \circ G$
(Theorem~2.3.1), $\widetilde{P}$ is order-invariant in $T \cap V$, with the
same constant order as~$P$ on~$S$ (by the remark following the definition
of order-invariance in Section~3.2 of the thesis).

\subsubsection{Complexification and order-invariance on
\texorpdfstring{$T^*$}{T*} (unchanged in method)}
\label{sssec:complexification}

\textbf{[UNCHANGED in method --- the argument applies to $\widetilde{P}$
instead of the discriminant $E$]}

Extend $g(y, y_r)$ to a pseudopolynomial $g(z, z_r)$ in a polydisc
$\Delta_1 \subseteq \C^{r-1}$ by extending each coefficient via its
convergent power series (Theorem~2.1.2).

Similarly, $\widetilde{P}(y)$ extends to a holomorphic function
$\widetilde{P}(z)$ in~$\Delta_1$.  The power series expansion of
$\widetilde{P}(y)$ about~$0$ has real coefficients; substituting complex
variables gives $\widetilde{P}(z)$.

Let $T^* = \{ z \in \C^{r-1} : z_{s+1} = \cdots = z_{r-1} = 0 \}$ be the
complexification of~$T$.  The order-invariance argument from the thesis
(pp.~25--26) applies verbatim to~$\widetilde{P}$:

Let $\mu = \ord_0 \widetilde{P}$.  Every partial derivative of
$\widetilde{P}(z)$ of order $< \mu$ vanishes throughout
$T^* \cap \Delta_1$ (because the corresponding real partial derivative
vanishes on $T \cap V$ by order-invariance, and the power series with
real coefficients extends this to the complex domain).  Some partial
derivative of order~$\mu$ does not vanish at~$0$.  After refining
$\Delta_1$, $\widetilde{P}(z)$ is order-invariant in $T^* \cap \Delta_1$
with constant order~$\mu$.

\subsubsection{Weierstrass preparation (unchanged)}
\label{sssec:weierstrass}

\textbf{[UNCHANGED]}

By the Weierstrass preparation theorem (Theorem~2.1.10), after refining
$\Delta_1$ to $\Delta_2 \subseteq \Delta_1$, there is a unit
$u(z, z_r)$ (holomorphic and non-vanishing in
$\Delta' = \Delta_2 \times \Delta(0; \epsilon)$) and a Weierstrass
polynomial
$h(z, z_r) = z_r^m + a_1(z) z_r^{m-1} + \cdots + a_m(z)$
in~$\Delta_2$ such that
\begin{equation}
  g(z, z_r) = u(z, z_r) \, h(z, z_r)
  \tag{3.3.7}
\end{equation}
for all $(z, z_r) \in \Delta'$, and for each fixed $z \in \Delta_2$,
all roots of $h(z, z_r)$ lie in $\Delta(0; \epsilon)$.

\subsubsection{Recovering order-invariance of \texorpdfstring{$\disc(h)$}{disc(h)}
  (changed --- key new step)}
\label{sssec:recover_discr}

\textbf{[CHANGED --- this is the main new argument]}

In the original proof, the discriminant $E(z) = \disc(g)$ is decomposed as
$E(z) = Q(z) \cdot F(z)$ where $F = \disc(h)$ and $Q$ involves
$\disc(u)$ and $\res(u, h)$, using Theorem~2.3.3.  We replace this with
the following argument.

\medskip

Since $u$ is a unit (holomorphic and non-vanishing), the ideals
$\langle g, g' \rangle$ and $\langle h, h' \rangle$ coincide in the ring
of holomorphic functions on~$\Delta'$:
\begin{align*}
  g &= uh, \quad g' = u'h + uh' \\
  \Longrightarrow\quad
  h &= u^{-1} g \in \langle g, g' \rangle, \quad
  h' = u^{-1}(g' - u' u^{-1} g) \in \langle g, g' \rangle \\
  \Longrightarrow\quad
  &\langle h, h' \rangle \subseteq \langle g, g' \rangle.
\end{align*}
The reverse inclusion is immediate from $g = uh$ and $g' = u'h + uh'$.

Since
$\widetilde{P}(z) \in \langle g, g' \rangle \cap \mathcal{O}(\Delta_2)
= \langle h, h' \rangle \cap \mathcal{O}(\Delta_2)$,
Corollary~\ref{cor:elim_weierstrass} gives
\begin{equation}
  \widetilde{P}(z) = \disc(h)(z) \cdot Q(z)
  \label{eq:P_disc_Q}
\end{equation}
for some holomorphic $Q \in \mathcal{O}(\Delta_2)$.

We verify that neither factor is identically zero:
\begin{itemize}
\item $\disc(h) \not\equiv 0$ because $f$ is squarefree, so $g$ is
  squarefree as a pseudopolynomial (the coordinate change preserves
  coprimality of $g$ and $g'$), and $h$ divides $g$ with unit cofactor,
  so $h$ is also squarefree.
\item $Q \not\equiv 0$ because $\widetilde{P} \not\equiv 0$ and
  $\disc(h) \not\equiv 0$, and $\mathcal{O}(\Delta_2)$ is an integral
  domain.
\end{itemize}

Now apply the order additivity lemma (Lemma~\ref{lem:order_additivity}) to
the restriction of~\eqref{eq:P_disc_Q} to $T^* \cap \Delta_2$:
\begin{itemize}
\item $\widetilde{P}|_{T^*}$ has constant vanishing order $\mu$ on
  $T^* \cap \Delta_2$ (established in \S\ref{sssec:complexification}).
\item $\disc(h)|_{T^*}$ and $Q|_{T^*}$ are holomorphic, not identically
  zero, on the connected open set $T^* \cap \Delta_2 \cong \Delta^{(s)}$.
\item Their product equals $\widetilde{P}|_{T^*}$, which has constant
  order.
\end{itemize}
By Lemma~\ref{lem:order_additivity}, $\disc(h)|_{T^*}$ has constant
vanishing order on $T^* \cap \Delta_2$.  Denote this constant order by~$r$.
Then $F(z) := \disc(h)(z) = z_{n-1}^r \cdot N(z)$ (after possibly
refining $\Delta_2$ and noting the power-series structure), where $N$ is
holomorphic and non-vanishing --- exactly the hypothesis required by
Zariski's theorem.

\subsubsection{Application of Zariski's theorem (unchanged)}
\label{sssec:zariski}

\textbf{[UNCHANGED]}

Since $\disc(h)$ is order-invariant in $T^* \cap \Delta_2$, Zariski's
1975 theorem (Theorem~4.1.1) applies to~$h$:
there exists a polydisc $\Delta_3 \subseteq \Delta_2$ and a holomorphic
function $\psi(z_1, \ldots, z_s)$ from $\Delta_3^{(s)}$ into
$\Delta(0; \epsilon)$ such that for all
$z = (z_1, \ldots, z_s, 0, \ldots, 0) \in T^* \cap \Delta_3$ and all
$\alpha \in \Delta(0; \epsilon)$,
\[
  h(z, \alpha) = 0 \quad\text{if and only if}\quad
  \alpha = \psi(z_1, \ldots, z_s),
\]
and $h$ is order-invariant in the set
$G^* = \{ (z, z_r) : z \in T^* \cap \Delta_3,\;
z_r = \psi(z_1, \ldots, z_s) \}$.

\subsubsection{Return to the real domain (unchanged)}
\label{sssec:return_real}

\textbf{[UNCHANGED]}

The argument on pp.~27--28 of the thesis applies verbatim.  In summary:
\begin{enumerate}
\item The power series of $\psi$ has real coefficients (by the
  Schwarz reflection argument), so $\eta := \psi|_{B^{(s)}}$ is
  real-valued and analytic.
\item For $y \in T \cap B$ and $\alpha \in (-\epsilon, +\epsilon)$:
  $g(y, \alpha) = 0 \iff h(y, \alpha) = 0 \iff \alpha = \eta(y)$
  (using $u \ne 0$).
\item Order-invariance of $g$ in the graph of~$\eta$ follows from
  order-invariance of $h$ in $G^*$ (since $g = uh$ and $u \ne 0$).
\item Detransforming via~$\Phi$: set $N = \Phi^{-1}(B)$ and
  $\theta(x) = (\eta \circ \phi)(x)$, where $\phi$ is the chart for~$S$
  corresponding to~$\Phi$.  This gives the desired analytic root function
  and order-invariance.  \qed
\end{enumerate}

%% =========================================================================
\section{The Generalized Projection Theorem (Theorem~3.2.3')}
\label{sec:projection}
%% =========================================================================

\textbf{[CHANGED --- proved from 3.2.1' via the Nullstellensatz]}

\begin{theorem}[Generalized Projection Theorem, Theorem 3.2.3']
\label{thm:gen_projection}
Let $A$ be a finite squarefree basis of $r$-variate integral polynomials,
$r \ge 2$, and let $S$ be a connected submanifold of\/ $\R^{r-1}$.
Suppose that each element of~$A$ is not identically zero on~$S$, and
each element of $P'(A)$ (Definition~\ref{def:gen_projection}) is
order-invariant in~$S$.  Then:
\begin{enumerate}
\item Each element of~$A$ is degree-invariant and analytically delineable
  on~$S$.
\item The sections of~$A$ over~$S$ are pairwise disjoint.
\item Each element of~$A$ is order-invariant in every section of~$A$
  over~$S$.
\end{enumerate}
\end{theorem}

\begin{proof}
Let $F_1, \ldots, F_n$ be the nonconstant elements of~$A$, and set
$f = \prod_{i=1}^{n} F_i$.  Then $f$ is squarefree (the $F_i$ are
squarefree and pairwise coprime, so
Lemma~\ref{lem:prod_squarefree} applies).

By hypothesis, we have nonzero order-invariant elements
\[
  d_i \in \langle F_i, F_i' \rangle \cap \R[x], \qquad
  r_{ij} \in \langle F_i, F_j \rangle \cap \R[x].
\]
Set $R = \prod_i d_i \cdot \prod_{i < j} r_{ij}$.

\begin{claim}
Some positive power $R^N$ belongs to
$\langle f, f' \rangle \cap \R[x]$.
\end{claim}

\begin{proof}[Proof of Claim]
Over $\bar{K}$ (algebraic closure), the repeated-root locus of~$f$
decomposes as
\[
  V(f, f') = \bigcup_i V(F_i, F_i') \cup \bigcup_{i < j} V(F_i, F_j).
\]
This holds because $f = \prod F_i$ with the $F_i$ pairwise coprime and
squarefree: a point $(a, t)$ is a repeated root of~$f$ iff either some
$F_i$ has a repeated root there, or two distinct factors vanish
simultaneously.

Since $d_i \in \langle F_i, F_i' \rangle \cap K[x]$, the polynomial $d_i$
vanishes on $\pi(V(F_i, F_i'))$.  Similarly $r_{ij}$ vanishes on
$\pi(V(F_i, F_j))$.  Therefore $R$ vanishes on
$\pi(V(f, f'))$.

By the Hilbert Nullstellensatz,
$R \in \sqrt{\langle f, f' \rangle \cap \bar{K}[x]}$,
so $R^N \in \langle f, f' \rangle \cap \bar{K}[x]$ for some $N \ge 1$.
Since $f$, $f'$, and $R$ all have integer coefficients, this membership
descends to $\R[x]$ (by flatness, or by a Gr\"obner basis computation
over~$\Q$).
\end{proof}

$R^N$ is a nonzero element of $\langle f, f' \rangle \cap \R[x]$.  It is
order-invariant on~$S$: each $d_i$ and $r_{ij}$ is order-invariant, so
$R$ is order-invariant by Lemma~\ref{lem:order_mul}, and
$\ord_a(R^N) = N \cdot \ord_a(R)$ is constant on~$S$.

Furthermore, $f$ is not identically zero on~$S$ and is degree-invariant
on~$S$ (the coefficients of each $F_i$ are order-invariant by hypothesis
on $\coeff(A)$).

Applying Theorem~\ref{thm:gen_lifting} (Theorem~3.2.1') to $f$ with
$P = R^N$: $f$ is analytically delineable on~$S$ and order-invariant in
each $f$-section over~$S$.  By Lemma~\ref{lem:order_mul}, each $F_i$ is
order-invariant in every $f$-section.  The sections of~$A$ over~$S$ are
pairwise disjoint (since $f$ is squarefree and delineable, distinct
coprime factors cannot share a root on any section).
\end{proof}

\begin{lemma}[used above]
\label{lem:prod_squarefree}
The product of squarefree pairwise coprime polynomials is squarefree.
\end{lemma}
\begin{proof}
Standard.
\end{proof}

%% =========================================================================
\section{Summary of changes from the thesis}
\label{sec:summary}
%% =========================================================================

\begin{center}
\begin{tabular}{|l|c|p{7cm}|}
\hline
\textbf{Item} & \textbf{Status} & \textbf{Notes} \\
\hline
Ch.~2 (all) & unchanged & Background material \\
\hline
Def.\ of $P(A)$ & \textbf{changed} & Generalized to $P'(A)$ with
  elimination ideal elements \\
\hline
Thm.~3.2.1 statement & \textbf{changed} & $P \in \langle f, f' \rangle \cap \R[x]$
  replaces discriminant~$D$ \\
\hline
Thm.~3.2.3 statement & \textbf{changed} & $P'(A)$ replaces $P(A)$ \\
\hline
Proof of 3.2.1, Assertion~1 from 2 & unchanged & \\
\hline
Proof of 3.2.1, Assertion~2, $s = r-1$ & \textbf{simplified} & $P \ne 0$
  directly rules out common roots \\
\hline
Proof of 3.2.1, Assertion~2, $1 \le s \le r-2$:
  coord.\ change & unchanged & \\
\hline
Proof of 3.2.1, $\widetilde{P}$ definition & \textbf{changed} &
  $\widetilde{P} = P \circ \Phi^{-1}$ instead of $E = \disc(g)$ \\
\hline
Proof of 3.2.1, complexification & unchanged in method &
  Same argument applied to $\widetilde{P}$ \\
\hline
Proof of 3.2.1, Weierstrass prep & unchanged & \\
\hline
Proof of 3.2.1, factorization step & \textbf{changed} &
  $\widetilde{P} = \disc(h) \cdot Q$ via
  Cor.~\ref{cor:elim_weierstrass} + order additivity
  (Lem.~\ref{lem:order_additivity}) replaces Thm.~2.3.3 \\
\hline
Proof of 3.2.1, Zariski application & unchanged & \\
\hline
Proof of 3.2.1, return to real & unchanged & \\
\hline
Proof of 3.2.3 & \textbf{changed} & Nullstellensatz argument replaces
  Thm.~2.3.3 product formula \\
\hline
Ch.~4 (Zariski theorem) & unchanged & Used as black box \\
\hline
\textbf{New:} Order additivity lemma & \textbf{new} &
  Lemma~\ref{lem:order_additivity} \\
\hline
\textbf{New:} Elim.\ ideal for Weierstrass & \textbf{new} &
  Cor.~\ref{cor:elim_weierstrass} (standard, citation-level) \\
\hline
\end{tabular}
\end{center}

%% =========================================================================
\section{Status relative to the Lean formalization}
\label{sec:lean_status}
%% =========================================================================

The Lean project already contains:
\begin{itemize}
\item \texttt{GeneralizedProjection.lean}: Theorem~3.2.3'
  (\texttt{mccallum\_3\_2\_3\_generalized}), fully proved from the
  generalized lifting theorem, which is currently an \textbf{axiom}
  (\texttt{lifting\_theorem\_generalized}).
\item The Nullstellensatz argument
  (\texttt{nullstellensatz\_elim}) is fully formalized.
\item All supporting lemmas (order additivity of products, squarefreeness
  of products, degree-invariance from coefficients, etc.) are formalized.
\end{itemize}

\noindent
\textbf{Remaining work to remove the axiom:}
\begin{enumerate}
\item Formalize the order additivity lemma
  (Lemma~\ref{lem:order_additivity}) for holomorphic functions.
  This requires Lean infrastructure for holomorphic functions on polydiscs
  and vanishing order in the complex-analytic setting.
\item Formalize Corollary~\ref{cor:elim_weierstrass} (elimination ideal
  for Weierstrass polynomials = discriminant).  This needs the resultant
  theory for monic polynomials over $\mathcal{O}(\Delta)$.
\item Formalize the proof of Theorem~3.2.1' (this document), using the
  above.  The bulk of the proof (coordinate changes, complexification,
  Weierstrass preparation, Zariski's theorem, return to real) is shared
  with the original theorem.
\end{enumerate}

\end{document}
